% Vorgaben Assignment aus Studienheft SQL03
% Formatvorgaben fuer den Text
% Umfang: 8 - 10 Seiten (inkl. Abbildungen und Tabellen, aber ohne Deckblatt, % Gliederung und Literaturverzeichnis, Eidesstattliche Erklaerung)
% Zeilenabstand: 1,5
% Schriftart: frei
% Schriftgrad: 12 pt
% Variablen, physikalische Groessen und Funktionszeichen werden kursiv gedruckt.
% Korrekturrand: links: 4,5 cm, rechts 2,0 cm, oben und unten jeweils 3,0 cm
% Deckblatt: (Adresse, AKAD-E-Mail-Adresse, Immatrikulationsnummer, Modul-
% bezeichnung, Thema, Datum, Felder für Korrektor)
% Gliederung (1 Seite)
% Literaturverzeichnis (3 - 5 Literaturquellen  z. B. Lehrbuecher, aktuelle Fachartikel recherchieren)
% Eidesstattliche Erklaerung (unterschrieben und fest eingebunden)
% Bearbeitungsdauer: 2 Monate


\documentclass[a4paper,12pt]{article}
\usepackage[ngerman]{babel}
\usepackage[nottoc]{tocbibind} % Anzeigen des Literaturverzeichnisses im TOC
\usepackage{epsfig}
\usepackage{times}
\usepackage{supertabular}
\usepackage{wrapfig}
\usepackage{multirow}
\usepackage[onehalfspacing]{setspace}
\usepackage{listings}
\usepackage{mathptmx}
\usepackage{geometry}
\usepackage{helvet}
\usepackage{courier}
\usepackage{setspace}
\usepackage{textcomp}
\usepackage[T1]{fontenc}
\usepackage[utf8]{inputenc}
\usepackage{fancyhdr}
\usepackage{float} % Notwendig fuer figure[h]
\usepackage[printonlyused]{acronym}

\newif\iflistoffigures
\newif\iflistoftables
\newif\ifacronym

%Titel
\newcommand*{\Titel}{SWE02 -- Sicherheitsaspekte eines Webshops} 

%Betreff
\newcommand*{\Betreff}{Assignment im Modul SWE02 -- \\ Softwareentwicklung \LaTeX{}} 

%Betreuer
\newcommand*{\Betreuer}{Betreuer: Prof. Kirchberg} 

%Vor- und Nachname
\newcommand*{\Name}{Stefan Waidele}

%Straße und Hausnummer
\newcommand*{\Strasse}{Ensisheimer Straße 2} 

%Plz und Ort
\newcommand*{\PlzOrt}{79395 Neuenburg am Rhein} 

%Immatrikulationsnummer
\newcommand*{\Immatrikulationsnummer}{102 81 71}

%Email 
\newcommand*{\Email}{Stefan.Waidele@AKAD.de} 


% Verzeichnisse (Wenn nicht benötigt, Zeile mit % auskommentieren oder löschen

%% Abbildungsverzeichnis 
%\listoffigurestrue
%% Tabellenverzeichnis
%\listoftablestrue
%% Abkürzungsverzeichnis
\acronymtrue


% Wittwen und Waisen verhindern
\clubpenalty10000
\widowpenalty10000
\displaywidowpenalty=10000

\usepackage[flushmargin,hang,ragged]{footmisc}
\usepackage{lmodern} %Type1-Schriftart für nicht-englische Texte
\usepackage{fancyhdr}

\usepackage[
	pdftitle={\Titel},
	pdfsubject={WEB04 -- Multichannel Marketing: Webshop, Filiale, Katalog},
	pdfauthor={Stefan Waidele},
	pdfkeywords={webshop, multichannel marketing, akad, web04, assignment, wirtschaftsinformatik}
	hyperfootnotes=false,
	colorlinks=true,
	linkcolor=black,
	urlcolor=black,
	citecolor=black
]{hyperref}

%\renewcommand{\familydefault}{\rmdefault}
\renewcommand{\bflabel}[1]{\normalfont{\normalsize{#1}}\hfill}

%% Definition for Codeschnipsel im Fließtext
\newcommand{\code}{\texttt}

%% Für Codeblöcke mit Syntax-Highlighting
%% http://www.ctan.org/tex-archive/macros/latex/contrib/minted/
\usepackage{minted}
\definecolor{bg}{rgb}{0.95,0.95,0.95}


\makeatother

\geometry{a4paper, left=45mm, right=20mm, top=30mm, bottom=30mm}
\pagenumbering{roman}

\begin{document}
	
\parskip=1em
\parindent=0cm

%description: Deckblatt in Deutsch
%% Basierend auf einer TeXnicCenter-Vorlage von Tino Weinkauf.
%% sowie der akad-vorlage von Daniel Falkner
%%%%%%%%%%%%%%%%%%%%%%%%%%%%%%%%%%%%%%%%%%%%%%%%%%%%%%%%%%%%%%

%%%%%%%%%%%%%%%%%%%%%%%%%%%%%%%%%%%%%%%%%%%%%%%%%%%%%%%%%%%%%
%% Deckblatt
%%%%%%%%%%%%%%%%%%%%%%%%%%%%%%%%%%%%%%%%%%%%%%%%%%%%%%%%%%%%%
%%
%% ACHTUNG: Sie benötigen ein Hauptdokument, um diese Datei
%%          benutzen zu können. Verwenden Sie im Hauptdokument
%%          den Befehl "\input{dateiname}", um diese
%%          Datei einzubinden.
%%

\begin{titlepage}

%\thispagestyle{empty}

\vspace{5cm}

\Name \\ 
\Strasse \\ 
\PlzOrt\\ 
\href{mailto:\Email}{\Email}

AKAD University\\
Immatrikulationsnummer: \Immatrikulationsnummer

\vfill

\begin{tabbing}
Modul WEB04 --- \=Fallstudie Online--Shop\\ 
                \>Assignment  
\end{tabbing}
\LARGE
\textsc{Multichannel--Marketing\\Webshop --- Filialen --- Katalog}\\

\vfill

\normalsize

Betreuer: Mark Harwardt, M.Sc.

\today %%Datum der Abgabe - am besten selbst reinschreiben.

\vfill

\includegraphics[width=3cm]{akad_logo.png}

AKAD University

\end{titlepage}


%\includegraphics[scale=0.35]{akad_logo.png}

%\clearpage
\normalsize


\normalsize

\begin{spacing}{1.0} % Verzeichnisse werden mit einzeiligem Abstand gesetzt
\parskip=0em
\newpage

% Inhaltsverzeichnis
\setcounter{tocdepth}{2}
\tableofcontents 
\newpage

% Abbildungsverzeichnis
%\iflistoffigures
\listoffigures 
\newpage
%\fi

% Tabellenverzeichnis
\iflistoftables
\listoftables 
\newpage
\fi

% Abkürzungsverzeichnis
\ifacronym
\section*{Abkürzungsverzeichnis}
\addcontentsline{toc}{section}{Abkürzungsverzeichnis} 
\begin{acronym}[\hspace{5cm}]
	\acro{CAPTCHA}{Completely Automated Public Tourint Test to Tell Computers and Humans Apart}
	\acro{CSS}{Cascading Style Sheets}
	\acro{DoS}{Denial of Service}
	\acro{DDoS}{Distributed Denial of Service}
	\acro{ERP}{Enterprise Resource Planning}
	\acro{HTML}{Hypertext Markup Language}
	\acro{HTTP}{Hypertext Tranfsfer Protocol}
	\acro{HTTPS}{Hypertext Tranfsfer Protocol Secure}
	\acro{OWASP}{Open Web Application Security Project}
	\acro{PDO}{PHP Data Objects}
	\acro{SCM}{Supply Chain Management}
	\acro{SQL}{Structured Query Language}
	\acro{WWW}{World Wide Web}
	\acro{XSS}{Cross Site Scripting}
	\acro{URL}{Unified Resource Locator, Internetadresse}

%%
%	
%	\acro{}{}
%	\acro{}{}
%	\acro{}{}
%	\acro{}{}
%	\acro{}{}
%	\acro{}{}
%	\acro{}{}
%	\acro{}{}
	\acro{Test}{Wird nicht im Text verwendet und taucht auch nicht im Verzeichnis auf}
\end{acronym}

\fi

\parskip=1em
\end{spacing} 

\clearpage

\newcounter{romanPagenumber} 
\setcounter{romanPagenumber}{\value{page}} % Roemische Seitenanzahl speichern.


\pagestyle{fancy}
\fancyhead{}
\fancyhead[LO,RE]{\textsc{\Titel}}
\fancyhead[RO,LE]{\thepage}
\fancyfoot[CO,CE]{}
\setlength{\headheight}{15pt}

% \nocite{*} 

\pagenumbering{arabic}

\begin{spacing}{1.5} % Zeilenabstand: 1,5 fuer den Textteil

\section{Einleitung}
\label{sec:einleitung}


\subsection{Ziele dieser Arbeit}

\todo{Aufgabenstellung: Ausgangssituation: Sie sind Leiter Marketing bei Manufactum, einem Versandhändler für qualitativ hochwertige Produkte. Neben dem klassischen Versandhandel von Katalogen gibt es mittlerweile auch einen Online-Shop sowie mehrere Filialen für den stationären Handel. Als Leiter Marketing liegt der Bereich E-Commerce in Ihrer Verantwortung.  

Hinweis zu den Aufgaben: Bitte beachten Sie, dass bei allen Aufgaben nötig ist, ein theoretisches Fundament zu schaffen.

Aufgabe E – Multichannel: Sie wollen Webshop, Filialen und Print-Kataloge sinnvoll miteinander „verdrahten“. Welche Möglichkeiten sehen Sie hier?
}

\subsection{Aufbau der Arbeit}

Zunächst werden im Kapitel~\myref{sec:grundlagen} erörtert, ...

Im Kapitel~\myref{sec:hauptteil} wird gezeigt, ...

\subsection{Abgrenzung}

\todo{Keine Unterscheidung zwischen Multichannel--Marketing, --Commerce, --Management, --Retailing, --Whatever}


\section{Grundlagen}
\label{sec:grundlagen}


\subsection{Definition Retailing}

\begin{quote}\textbf{retail:} „The sale of goods to the public in relatively small quantities for use or consumption rather than for resale.“\footnote{\cite{oxford:retail}}\end{quote}

Bei Manufactum handelt es sich nach obiger Definition um einen Akteur auf der Einzelhandelsebene: Die Waren werden in überschaubarer Menge an Endkunden verkauft. Die Vertriebswege Katalogversand, Ladengeschäft und Onlineshop sind ebenfalls im Einklang mit der von Wirz beschriebenen Einteilung\footnote{vgl. \cite{wirtz}, S.33}.

International spricht man bei der Kombination von Internetshop, Ladengeschäften und Katalogen von „Clicks, Bricks and Sheets“\footnote{vgl. \cite{klein}}.

\subsection{Definition Multi--Channel--Marketing}

In dieser Arbeit wird der Begriff Multichannel Marketing gemäß folgender Definitionen von Wirtz bzw. Emrich (siehe nächster Abschnitt) verwendet:

\begin{quote}„Unter \textbf{Multi--Channel--Marketing} versteht man den Prozess der Planung, DUrchführung und Kontrolle aller Marketingaktivitäten in einem Mehrkanalsystem. Dadurch sollen durch eine dauerhafte Befriedigung der Kundenbedürfnisse die Unternehmensziele verwirklicht werden.“\footnote{\cite{wirtz}, S21}\end{quote} 

\subsection{Definition Mehrkanalsystem}

Der Begriff „Mehrkanalsystem“ wird von Emrich in der oben zitierten Definition konkretisiert:

\begin{quote}„Mult--Channel--Marketing ist die Nutzung mehrerer multifunktional vernetzter Kanäle sowohl für Kommunitkations als auch für Vertrieb von Produkten/Dienstleistungen eines Anbieters an organisationale Kunden bzw. Endverbraucher; es enthält mindestens zwei eigenständige unterschiedliche Kanäle für markierte Leistungsbündel mit einem Sortimentszusammenhang, für die ein kanalspezifischer Marketing--Mix bestehen kann und die in ein ganzheitliches Channel--Konzept integriert bzw. mit diesem kombiniert sind.“\footnote{\cite{emrich}, S. 8}\end{quote} 

Aufgrund der Entwicklungen im Bereich Mobilfunk bzw. mobiles Internet wird dieser Kanal als besonders wichtig für Multichannel--Strategien angesehen\footnote{vgl. \cite{mohlenbruch}}.

\subsection{Verknüpfungsmöglichkeiten zwischen Kanälen}

Bei den drei Verkaufskanälen Webshop, Ladengeschäft und Katalogbesteht die Notwendigkeit die folgenden Verknüpfungen herzustellen:

\begin{itemize}
\item{\textbf{Katalog --- Ladengeschäft}\\
Der Hinweis auf persönliche Beratung in den Ladengeschäften ist bereits im Katalog vorhanden\footnote{siehe \cite{manufactum}, S. 403}. Ebenfalls sind die Öffnungszeiten, Adressen und Kontaktmöglichkeiten abgedruckt. Die einzelnen Standorte werden ausführlich und attraktiv vorgestellt.
}
\item{\textbf{Katalog --- Webshop}\\
Wie auf die Ladengeschäfte wird im Katalog ebenfalls auf die Angebote im Internet hingewiesen\footnote{siehe \cite{manufactum}, S. 396}. Auch der Internetshop wird vorgestellt. Die Verweise ins Internet werden durch die jeweilige \ac{URL} getätigt.\\
Eine weitere Möglichkeit der Übermittlung von URLs besteht in der Nutzung von QR--Codes, welche von Mobiltelefonen. Auch hiervon wird im Katalog Gebrauch gemacht, allerdings nur vereinzelt.
}
\item{\textbf{Ladengeschäft --- Katalog}\\
Durch die Auslage von Katalogen und Prospekten in den Ladengeschäften wird diese Verknüpfung hergestellt.
}
\item{\textbf{Ladengeschäft --- Webshop}\\
Kunden im Ladengeschäft können durch die Anbringung von \ac{URL}s oder QR--Codes auf Werbemitteln wie z.B. Preisschildern auf den Webshop aufmerksam gemacht werden.
}
\item{\textbf{Webshop --- Ladengeschäft}\\
Im Webshop sind außer den mit Print vergleichbaren Hinweisen noch spezielle Informationen bzgl. der Verfügbarkeit einzelner Produkte in den Geschäften denkbar. 
}
\item{\textbf{Webshop --- Katalog}\\
Im Webshop bzw. auf den Webseiten von Manufactum wird auf Bestell-- und Downloadmöglichkeiten der diveren Katalogen hingewiesen.
}
\end{itemize}
\section{Hauptteil}
\label{sec:hauptteil}

\subsection{Kurz--URLs}

Die zunächst einfachst Verbindung von Gedrucktem zu Internetseiten bzw. zum Webshop besteht darin, dem Kunden eine \ac{URL}, also eine Internetadresse\footnote{Obwohl der Begriff URL eine Internetadresse in sehr allgemeiner Form beschreibt, wird er in dieser Arbeit der Lesbarkeit halber synonym zu „WWW--Adresse“ verwendet.} zu nennen, die dieser dann am PC oder per Mobiltelefon aufruft.
Soll lediglich auf die Startseite des Webangebots hingewiesen werden, ist hierzu die gewöhnliche Adresse hier also \url{www.Manufactum.de}, anzugeben. 

Problematischer wird es, wenn Links zu spezifischen Produktseiten in Webshop abgedruckt werden sollen. Die hier entstehenden \ac{URL}s sind i.d.R. eher lang und somit umständlich für den Kunden:\\ \url{http://www.manufactum.de/manufactum-ziehharmonikaboerse-klein-p759940/?a=64182}

Daher empfielt sich der Einsatz von sogenanten Kurz--URLs. Hier wird eine neue, deutlich kürzere \ac{URL} erzeugt, die den Besucher beim Aufruf auf die ursprüngliche Adresse umleitet. Technisch wird dies dadurch erreicht, dass der Hashwert einer langen Adresse an den Domainnamen des URL--Shorteners gehängt wird. Dadurch können viele lange Adressen auf jeweils
entsprechende Kurzlinks abgebidet werden: \url{http://goo.gl/QCJH55}\\
Dienste wie Ow.ly, Bit.ly, Goo.gl oder FixURL.de bieten entsprechende Dienste an, die per \ac{API} auch automatisiert genutzt werden können\footnote{siehe z.B. \url{http://dev.bitly.com/}}. Meist stellen diese Anbieter auch zumindest Aufrufstatistiken\footnote{z.B. Goo.gl} bis hin zu kompletten Analysedaten\footnote{z.B. Ow.ly durch das Social Media Management Tool Hootsuite} zur Vwerfügung.

Da der Kunde einer durch einen Dienst gekürzten \ac{URL} nicht mehr ansieht, wohin sie eigentlich führt, kann man durch die Verwendung eines eigenen, der Marke eindeutig zuordenbaren und dennoch kurzen Domainnamen verwenden. Hierdurch hat man dann volle Kontrolle über die erzeugten Links und Statistiken.\footnote{vgl. \cite{webmag} sowie \cite{gillen}}

\subsection{QR--Codes}

\ac{QR--Codes} 

\subsection{Verfügbarkeitsanzeige bzw. Bestellung zur Ansicht}
\subsection{Info--Terminal}


\section{Fazit \& Ausblick}

\subsection{Fazit}

Durch Umsetzung der Empfehlungen in Kapitel \myref{hempf} ist es Manufactum möglich, die vorhandenen Vertriebskanäle optimal zu verbinden. Dem Kunden bietet sich somit ein flexibles Einkaufserlebnis, das von Öffnungszeiten, räumlicher Distanz aber auch technischer Ausstattung unabhängig ist. 

Da dem Kunden eine medienübergreifende Auswahl an Einkaufsmöglichkeiten geboten wird, können die beliebtesten Wege durch Beobachtung der eigenen Daten ermittelt und ausgebaut werden. Das Risiko, den Schritt vom Katalogversender zum digital ausgerichteten Fachgeschäft zu verpassen wird minimiert. Neue Kundengruppen können angesprochen werden.

Wie anhand der Bildbeispiele zu sehen ist, handelt es sich bei den beschriebenen Möglichkeiten nicht um exotisches Neuland, sondern um etablierte Techniken, die bereits vielfach im Einsatz sind. Somit bedeutet die Umsetzung wohl weniger die Schaffung von Wettbewerbsvorteilen, sondern eher um die Verhinderung von entsprechenden Nachteilen gegenüber der Konkurrenz.

\subsection{Ausblick}

Die folgenden Themengebiete wurden in der vorliegenden Arbeit nicht behandelt, stellen jedoch Aspekte dar, die im Zusammenhang von Multi--Channel--Marketing untersucht werden sollten:

Durch die Einbeziehung von digitalen Techniken ist es möglich, Kundenwünsche besser zu erkennen und darauf zu reagieren. Dies kann in Form von Targetting\footnote{siehe z.B. \cite{omptarget}} oder allgemeiner durch Analyse und Optimierung der Customer Journey\footnote{siehe z.B. \cite{ompcj}} geschehen. Auch entstehen immer wieder neue Möglichkeiten, digital mit Kunden in Kontakt zu treten\footnote{Etwa NFC, Bluetooth, iBeacon, aber auch Messengerdienste wie Whatsapp oder Facebook--Messenger. Vgl. \cite{waidele}}. Diese sollten regelmäßig erhoben und evaluiert werden, um hier stets die von der Zielgruppe genutzten Kanäle zu bedienen.


Des weiteren ist auch eine fortlaufende Untersuchung der Rentabilität der einzelnen Maßnahmen notwendig. Dies gilt sowohl für die bereits genutzten Vertriebsöglichkeiten als auch für die noch nicht Implementierten.

\include{anhang}

\end{spacing}

\clearpage

% Literaturverzeichniss - Ab hier wieder Roemische Seitenzahlen
\pagestyle{plain}
\pagenumbering{roman}
\setcounter{page}{\theromanPagenumber}
\bibliographystyle{apalike}
\bibliography{literatur}
\onehalfspacing
\clearpage

\pagestyle{empty} 
\thispagestyle{empty}

\begin{center}
{\Large Eidesstattliche Erklärung}
\vspace*{4cm}\end{center}
\noindent
Ich versichere, dass ich das beiliegende Assignment selbstständig verfasst, keine anderen als die angegebenen Quellen und Hilfsmittel benutzt sowie alle wörtlich oder sinngemäß übernommenen Stellen in der Arbeit gekennzeichnet habe. 
\vspace{3cm}

\hspace{-0.8cm}
\rule[0.5ex]{6.5cm}{1pt}
\hspace{1.3cm}
\rule[0.5ex]{6.5cm}{1pt}
\\(Datum, Ort)
\hspace{6.3cm}
(Unterschrift)

\clearpage

%Messbox zur Druckkontrolle:
\newcommand{\Messbox}[2]{% Parameters: #1=Breite, #2=Hoehe
\setlength{\unitlength}{1.0mm}%
\begin{picture}(#1,#2)%
\linethickness{0.05mm}%
\put(0,0){\dashbox{0.2}(#1,#2)%
{\parbox{#1mm}{%
\centering\footnotesize 
%{\bf MESSBOX}\\ 
Breite $ = #1 {\rm\ mm}$\\
H\"ohe $ = #2 {\rm\ mm}$
}}}\end{picture}
}

\begin{center}
\textbf{--- Druckgröße kontrollieren! ---}
\\
\Messbox{100}{50} % Angabe der Breite/Hoehe in mm
\\
\textbf{--- Diese Seite nach dem Druck entfernen! ---}
\end{center}


\end{document}

